\documentclass{article}
\usepackage{graphicx}
\usepackage[brazilian]{babel}
\usepackage[latin1]{inputenc}
\usepackage{amssymb}
\usepackage{amsmath}
\usepackage{indentfirst}

\usepackage{xcolor}
\definecolor{verde}{rgb}{0,0.5,0}

\usepackage{listings}
\lstset{
  language=c,
  basicstyle=\ttfamily\small, 
  keywordstyle=\color{blue}, 
  stringstyle=\color{verde}, 
  commentstyle=\color{red}, 
  extendedchars=true, 
  showspaces=false, 
  showstringspaces=false, 
  numbers=left,
  numberstyle=\tiny,
  breaklines=true, 
  backgroundcolor=\color{green!10},
  breakautoindent=true, 
  captionpos=b,
  xleftmargin=0pt,
}

\begin{document}

\title{Lista de Exerc�cios 5 - Super Lista}
\author{PUC Minas - Engenharia de Computa��o \\ Algoritmos e Estrutura de Dados II \\ Prof. Kleber Jacques F. de Souza}
\date{}
\maketitle

%===================================================================================================

\section*{Instru��es}

\begin{itemize}
\item Esta � uma lista de refor�o e vale 2 pontos
\item No dia 15/04 deve ser entregue a vers�o manuscrita do c�digo solicitado (1 ponto)
\item No dia 21/04 deve ser entregue uma vers�o do c�digo fonte do c�digo solicitado implementado (1 ponto)
\end{itemize}

%===================================================================================================

\section*{Super Lista}

Desenvolva um Tipo Abstrato de Dado (TAD) chamado \textbf{Super Lista} que execute todas as opera��es abaixo:

\begin{enumerate}
\item Inserir um dado no come�o da lista
\item Inserir um dado no fim da lista
\item Inserir um dado na posi��o $p$ da lista
\item Verificar se a lista est� vazia
\item Retornar o tamanho da lista
\item Imprimir os dados da lista
\item Imprimir os dados da lista em ordem crescente
\item Remover um dado do come�o da lista
\item Remover um dado do fim da lista
\item Remover um dado da posi��o $p$ da lista
\item Pesquisar um elemento $e$ est� na lista
\item Verificar quantas vezes um elemento $e$ se repete na lista

\end{enumerate}

\textbf{A sua implementa��o deve ser feita utilizando aloca��o din�mica de mem�ria e ponteiros!}


%===================================================================================================

\end{document}